\BourbakiExercisesSection

\begin{enumerate}
\def\labelenumi{\arabic{enumi}.}
\tightlist
\item
  设 E 是一个集合， \(A \subset E \times E\) 与 \((x, y) \mapsto x \top y\) , \((x, y) \in A\) 是 E 上的一个未必处处有定义的内部合成律。给定 E 的两个子集 X 与 Y，令 \(X \top Y\) 表示 E 的形如 \(x \top y\) （其中 \(x \in X\) , \(y \in Y\) 且 \((x, y) \in A\) ）的元素所成的集合。这样就在 \(\mathfrak{P}(E)\) 上定义了一个处处有定义的合成律。
\end{enumerate}

若 \((\mathbf{X}_{\alpha})_{\alpha \in \mathbf{A}}\) 与 \((\mathbf{Y}_{\beta})_{\beta \in B}\) 是 E 的子集的任意两个族，则

\[
\left(\bigcup_ {\alpha \in A} X _ {\alpha}\right) \top \left(\bigcup_ {\beta \in B} Y _ {\beta}\right) = \bigcup_ {(\alpha , \beta) \in A \times B} (X _ {\alpha} \top Y _ {\beta}).
\]

\begin{enumerate}
\def\labelenumi{\arabic{enumi}.}
\setcounter{enumi}{1}
\item
  设 \(\top\) 是集合 E 上未必处处有定义的一个合成律。设 E' 为 \(\mathfrak{P}(E)\) 的由诸集合 \(\{x\}\) （其中 \(x\) 遍历 E）与 E 的空子集 \(\varnothing\) 所组成的子集；证明 E' 是 \(\mathfrak{P}(E)\) 在律 \((X,Y)\mapsto X\top Y\) （习题1）下的稳定子集；由此推出，若 \(\overline{E}\) 表示通过向 E 添加（《集合论》II，§ 4，第 8 号）一个元素 \(\omega\) 而得到的集合，则律 \(\top\) 可以延拓到 \(\overline{E} \times \overline{E}\) 上，使得律 \(\top\) 与该延拓律在 E 上诱导的律相同。
\item
  设 \(\top\) 是 E 上并非处处有定义的一个合成律。
\end{enumerate}

\begin{enumerate}
\def\labelenumi{(\alph{enumi})}
\item
  对于律 \((\mathbf{X}, \mathbf{Y}) \mapsto \mathbf{X} \top \mathbf{Y}\) （习题1）—— \(\mathbf{E}\) 的子集之间的律——要使它是结合的，必要且充分的是：对所有 \(x \in \mathbb{E}\) , \(y \in \mathbb{E}\) , \(z \in \mathbb{E}\) ，若公式 \((x \top y) \top z = x \top (y \top z)\) 的两侧中有一侧有定义，则另一侧也有定义且与之相等（利用习题1证明该条件是充分的）。
\item
  若此条件成立，证明第 3 节定理 1 可推广如下：若公式 (3) 的两侧中有一侧有定义，则另一侧有定义且与之相等。
\end{enumerate}

\begin{enumerate}
\def\labelenumi{\arabic{enumi}.}
\setcounter{enumi}{3}
\tightlist
\item
  \begin{enumerate}
  \def\labelenumii{(\alph{enumii})}
  \tightlist
  \item
    给定集合 E，设 \(\Phi\) 为 E 的任意子集到 E 中的映射所成的集合；若 f、g、h 是 \(\Phi\) 的三个元素，证明若复合 \((f \circ g) \circ h\) 有定义，则 \(f \circ (g \circ h)\) 也有定义，但反之不成立；若这两个复合都有定义，则它们相等。
  \end{enumerate}
\end{enumerate}

\begin{enumerate}
\def\labelenumi{(\alph{enumi})}
\setcounter{enumi}{1}
\tightlist
\item
  设 \(\mathfrak{F}\) 是 \(E\) 的一族非空子集，其中任意两个都无公共元素，并令 \(\Psi\) 为 \(\Phi\) 的由 \(\mathfrak{F}\) 中一个集合到 \(\mathfrak{F}\) 中一个集合上的双射映射所组成的子集。证明：在由律 \(f \circ g\) 在 \(\Psi\) 上诱导的律下，习题3(a)的条件成立。
\end{enumerate}

\begin{enumerate}
\def\labelenumi{\arabic{enumi}.}
\setcounter{enumi}{4}
\item
  证明自然数的三元组 \((m, n, p)\) （ \(\neq 0\) ）中，使得 \((m^n)^p = m^{n^p}\) 的只有： \((1, n, p)\) ，其中 \(n\) 与 \(p\) 任意； \((m, n, 1)\) ，其中 \(m\) 与 \(n\) 任意；以及 \((m, 2, 2)\) ，其中 \(m\) 任意。
\item
  设 \(\top\) 是集合 E 上的一个合成律；设 A 为 E 的由这样的元素 \(x\) 组成的子集：满足 \(x \top (y \top z) = (x \top y) \top z\) （对一切 \(y \in E\) , \(z \in E\) ）；证明 A 是 E 的稳定子集，且 \(\top\) 在 A 上诱导的律是结合的。
\item
  若 \(\top\) 是 \(\mathbf{E}\) 上的结合律， \(a\) 与 \(b\) 是 \(\mathbf{E}\) 的两个元素，则集合 \(\{a\} \top \mathbf{E}, \mathbf{E} \top \{b\}, \{a\} \top \mathbf{E} \top \{b\}, \mathbf{E} \top \{a\} \top \mathbf{E}\) 是 \(\mathbf{E}\) 在律 \(\top\) 下的稳定子集。
\item
  设 \(\top\) 是 E 上的一个结合律， \(a\) 是 E 的一个元素；对一切 \(x \in \mathrm{E}\) , \(y \in \mathrm{E}\) ，我们记 \(x \perp y = x \top a \top y\) ：证明律 \(\bot\) 是结合的。
\item
  在集合 E 上，映射 \((x,y)\mapsto x\) 与 \((x,y)\mapsto y\) 是互为相反的结合律。
\item
  设 X 与 Y 是集合 E 的任意两个子集，令 \(X \top Y = X \cup Y\) 若 \(X \cap Y = \varnothing\) ，令 \(X \top Y = E\) 若 \(X \cap Y \neq \varnothing\) ；证明这样在 \(\mathfrak{P}(E)\) 上定义的合成律是结合的且交换的。
\item
  设 \(\top\) 是 \(\mathbf{E}\) 上的一个结合律， \(\mathbf{A}\) 与 \(\mathbf{B}\) 是 \(\mathbf{E}\) 的在该律下稳定的两个子集；证明，若 \(\mathbf{B} \top \mathbf{A} \subset \mathbf{A} \top \mathbf{B}\) ，则 \(\mathbf{A} \top \mathbf{B}\) 是 \(\mathbf{E}\) 的稳定子集。
\item
  仅有的不同的自然数 \(\neq 0\) ，它们在律 \((x,y)\mapsto x^{y}\) 下可交换的，是 2 和 4。
\item
  证明在律 \((\mathbf{X}, \mathbf{Y}) \mapsto \mathbf{X} \circ \mathbf{Y}\) （ \(E \times E\) 的子集之间）下，中心是由 \(\varnothing\) 与对角线组成的集合。
\item
  证明在 E 到 E 的映射之间的合成律 \(f \circ g\) 下，中心由恒等映射组成。
\item
  律 \(\top\) （在集合 \(E\) 上）称为幂等的，如果 \(E\) 的所有元素在该律下都是幂等的（第 4 号），即若 \(x \top x = x\) 对一切 \(x \in E\) 成立。证明，若律 \(\top\) 在 \(E\) 上是结合的、交换的且幂等的，则关系 \(x \top y = y\) 是 \(E\) 上的一个序关系；若记 \(x \leqslant y\) ，则任意两个元素 \(x, y\) （属于 \(E\) ）都有一个最小上界（在该序关系下）等于 \(x \top y\) 。试求其逆。
\item
  设 \(\top\) 是集合 E 上的一个结合且交换的合成律。设 \(n\) 为整数 \(>0\) 。映射 \(x \mapsto \frac{n}{\top} x\) 是 E 到 E 中的一个态射。
\end{enumerate}

